% 张三慧《学习辅导与习题解答》PDF S 独立选答，卷1第3章；未取得V1主教材独立交叉核验。
% 仅录“思考题选答”中实际明确列出的原编号；其余不推断，exam_ready=false。
\section*{第3章\quad 动量与角动量：辅导书已列思考题独立解析}
\paragraph{思考题 3.3}
缓慢下拉时近似静平衡，下线张力为外拉力 \(F\)，上线同时承受 \(mg+F\)，故若相同破断强度，上线先断。若下拉冲量发生在很短时间，重球惯性使位移很小，力脉冲主要作用于下线；上线张力来不及同步增大，通常\emph{下线先断}。瞬间张力波传播、线弹性与拉力历时会影响真实破断情况，不可将其当任意拉法的绝对定理。

\paragraph{思考题 3.9}
设中心为O，质点质量m、半径R、角速度 \(\omega\)，关于O的 \(\boldsymbol L_O=mR^2\omega\hat z\) 不变，原因是向心合力对O力矩为零。对圆周上固定点A，力的作用线通常不穿过A，所以 \(\boldsymbol\tau_A\ne0\)，角动量一般随时间变。对轴上但不在O的固定点B，其 \(\boldsymbol L_B=\boldsymbol L_O+(\boldsymbol r_O-\boldsymbol r_B)\times m\boldsymbol v\)，第二项随速度方向转动，角动量矢量通常不守恒。只有选圆心时本题的完整角动量恒定。

\paragraph{思考题 3.10}
静止金核与α粒子间的库仑斥力指向两者连线，属于中心力，故关于\emph{固定核中心}的力矩 \(\boldsymbol\tau=\boldsymbol r\times\boldsymbol F=0\)，有 \(\boxed{\boldsymbol L_\alpha=m_\alpha\boldsymbol r\times\boldsymbol v=\mathrm{const}}\)。但α粒子受净电场力而改变动量方向，\(\dot{\boldsymbol p}_\alpha=\boldsymbol F\ne0\)，线动量不守恒。若考虑金核反冲，则整个\emph{两体系统}线动量及总角动量守恒，不能简单把α粒子相对于实验室固定核中心角动量视为严格恒定。
